\begin{abstract}
Severe brain injury, polyhandicap, and aphasia often lead to profound impairments in expressive communication, limiting caregivers' ability to detect internal emotional states such as pain, fear, discomfort, or positive affect. Electroencephalography-based brain-computer interfaces (EEG-BCIs) offer a promising avenue for restoring minimal communication channels in individuals unable to communicate verbally or through gestures. However, strong inter-patient variability in neural organization makes population-level models ineffective in this clinical context. 

We propose a patient-specific framework for emotion decoding using an 8-electrode Bitbrain Air EEG system combined with Self-Organizing Maps (SOMs) trained independently for each patient. In this technical validation study, we present results from a 4.7-minute recording session on a healthy female volunteer (age 42), demonstrating the complete pipeline from acquisition to classification. Despite the short duration, the system successfully extracted 279 analysis windows, identified 6 distinct emotional states from 9 caregiver annotations, and achieved 77.8\% classification accuracy (chance level = 16.7\%). Feature importance analysis revealed that connectivity variability (std\_abs\_corr = 0.138) and theta band power (0.125) were the most discriminative features, with frontal alpha asymmetry (FAA = 0.105) also contributing significantly.

The SOM successfully discovered structured representations in EEG feature space (quantization error = 2.34, topographic error = 0.038), with distinct clusters corresponding to different emotional states. Connectivity analysis revealed characteristic patterns for each state, and spectral profiles showed expected variations in theta and alpha bands. These results demonstrate that even with limited data, the patient-specific approach can extract meaningful neurophysiological correlates of emotional states.

This proof-of-concept establishes the technical foundation for practical affective BCIs in individuals with profound communication disorders. We present a comprehensive experimental plan for extending this work to eight polyhandicapped patients with aphasia in a future funded project, including detailed protocols for caregiver training, data acquisition during natural care activities, and longitudinal model validation.
\end{abstract}